La siguiente guía representa los tipos de datos que maneja Matlab:

\begin{itemize}
    \item Numeric: son vectores, matrices, flotantes, enteros, etcétera. El siguiente código es un ejemplo para un número random generado:
    
    \begin{lstlisting}
    %Datos numericos
    r = rand;
    >>> 0.157613081677548
    \end{lstlisting}

    \item Text (strings): todo lo que se encuentra encerrado entre corchetes o apóstrofos. Podemos transformar datos numéricos a strings mediante \texttt{num2str}, o viceversa \texttt{str2num}. Por ejemplo:
    
    \begin{lstlisting}
    %Datos text
    a = 42
    fprintf('El resultado es %s \n', num2str(a))
    \end{lstlisting}

    \item Structure: parecido a las estructuras que se definen en C++, donde defino propiedades al elemento. La única diferencia es que en MATLAB crea la variable de tipo struct. Un ejemplo es:
    
    \begin{lstlisting}
    %Datos structure
    auto(1).velocidad = 80;
    auto(1).nombre = 'espantomobil';
    auto(1).ruedas = 3;
    auto(2).ruedas = 4;

    ruedas_auto2 = getfield(auto(2), 'ruedas');

    personas(1).nombre = 'Fabrizio';
    personas(1).altura = 1.80;
    personas(2).nombre = 'Bautista';
    personas(2).altura = 1.65;

    function [alto] = masAlto(personas)
        if personas(1).altura > personas(2).altura
            alto = personas(1).nombre;
        else
            alto = personas(2).nombre;
        end
    end

    a = masAlto(personas);
    \end{lstlisting}

    \item Cell: igual que el structure pero accediendo a los datos mediante índices. Por ejemplo:

    \begin{lstlisting}
    %Datos cell
    bebidas{1, 1} = {'Coca', 15, 0};
    bebidas{2, 1} = {'Cerveza', 30, 5};
    bebidas{3, 1} = {'Fernet', 75, 40};
    bebidas{3, 1}{1}(1:4); %De esta forma cuento las primeras cuatro letras de la tercera bebida
    \end{lstlisting}

    \item Boolean: representan la lógica 'si' o 'no'. Permite correr ciertas partes de código si se cumple tal condición. Por ejemplo:
    
    \begin{lstlisting}
    %booleanos
    flag = true;
    if flag
        display('Se ejecuto esta parte del codigo')
    else
        display('El valor de flag es falso') 
    end
    \end{lstlisting}

    \item Function Handle: es un tipo de dato de MATLAB que representa una función. Por ejemplo:
    
    \begin{lstlisting}
    %function handle
    f1 = figure;
    f2 = figure;
    f3 = figure;

    figure(f1);
    set(f1, "Name", 'Primera figura', 'Position', [50 100 400 300]);
    set(f2, 'Color', [0.5, 0.1, 0.2])
    plot(1:10);
    \end{lstlisting}

    Este último llamamos a la figura 1 para trabajar con ella.
\end{itemize}